% L5: Gauss sums, Stickelberger element, Herbrand's theorem.

This chapter assembles the algebraic ingredients that give the Herbrand
implication: a non-trivial $\chi$-eigenspace in the $p$-part of the
class group forces $p \mid \Bgen{1}{\chi}$.  The analytic chapters below use
the same Gauss-sum normalization in the relative class-number formula.

\section{Gauss sums}

\begin{definition}\label{def:gauss-sum}
  \lean{gaussSum}\leanok
  \uses{def:dirichlet-character}
For a Dirichlet character $\chi$ modulo $p$ valued in $\C$, the \emph{Gauss
sum} is
\[
  \tau(\chi) \;=\; \sum_{a \in \Z/p\Z} \chi(a)\,e^{2\pi i a / p}
  \;\in\; \C,
\]
where the additive character is \(a \mapsto e^{2\pi i a/p}\).
\end{definition}

\begin{proposition}\label{prop:gauss-norm}
  \lean{BernoulliRegular.gaussSum_mul_gaussSum_inv_stdAddChar}\leanok
  \uses{def:gauss-sum}
For a non-trivial Dirichlet character $\chi$ modulo $p$,
\[
  \tau(\chi)\,\tau(\bar\chi) \;=\; \chi(-1)\,p,
  \qquad
  \abs{\tau(\chi)}^2 \;=\; p.
\]
\end{proposition}

\begin{proof}\leanok
Expand both Gauss sums and multiply:
\[
  \tau(\chi)\tau(\bar\chi)
  =
  \sum_{a,b}\chi(a)\bar\chi(b)e^{2\pi i(a+b)/p}.
\]
After the change of variables \(a=tb\) on the terms with \(b\ne0\), the
inner sum over \(b\) vanishes unless \(t=-1\).  The surviving contribution
is \(\chi(-1)p\).  Moreover
\(\overline{\tau(\chi)}=\chi(-1)\tau(\bar\chi)\), and therefore
\[
  \abs{\tau(\chi)}^2
  = \tau(\chi)\overline{\tau(\chi)}
  = \chi(-1)\tau(\chi)\tau(\bar\chi)
  = p.
\]
\end{proof}

The Gauss-sum sign (whether $\tau(\chi) = +\sqrt{p}$ or $-\sqrt{p}$ for
a real quadratic character) is fixed by the classical quadratic Gauss-sum
evaluation and is used in the $L(1,\chi)$ formula of
\cref{ch:lvalue-negative}.

\section{The Stickelberger element}

\begin{definition}\label{def:stickelberger}
  \lean{BernoulliRegular.stickelbergerElement}\leanok
  \uses{def:dirichlet-character}
The \emph{Stickelberger element} is
\[
  \theta \;=\; \sum_{a \in (\Z/p\Z)^\times}
                \Bigl\langle \frac{a}{p} \Bigr\rangle\,\sigma_a^{-1}
              \;\in\; \Q\bigl[(\Z/p\Z)^\times\bigr],
\]
where \(\langle x \rangle = x-\lfloor x\rfloor\) denotes the fractional
part, and \(\sigma_a\) is the automorphism sending
\(\zeta_p\mapsto\zeta_p^a\).  Multiplying by the appropriate denominator gives
the usual integral Stickelberger ideal.
\end{definition}

\begin{theorem}[Stickelberger's theorem]\label{thm:stickelberger}
  \lean{BernoulliRegular.stickelbergerCharacterCoefficientGroupRingTarget_annihilates_primeClass}
  \leanok
  \uses{def:stickelberger, def:gauss-sum}
The integral Stickelberger ideal $I_\theta \subset \Z[(\Z/p\Z)^\times]$
annihilates the class group of $\Qzetap$: for any prime $\mathfrak l \nmid p$
and any $\beta \in I_\theta$, the ideal $\mathfrak l^\beta$ is principal.
\end{theorem}

\begin{proof}\leanok
One first proves the ideal factorisation of Gauss sums:
for a non-trivial character \(\chi\), the principal ideal generated by
\(\tau(\chi)\) has factorisation prescribed by the Stickelberger exponent.
Since a principal ideal represents the trivial class, the corresponding
group-ring exponent annihilates the class of every prime above a rational
prime \(\ell\ne p\).  Such prime classes generate the ideal class group, so
every element of the integral Stickelberger ideal annihilates the full class
group.
\end{proof}

\section{Herbrand's theorem}

Herbrand's theorem converts Stickelberger annihilation into a divisibility
statement for generalised Bernoulli numbers:

\begin{theorem}[Herbrand bridge]\label{thm:herbrand}
  \lean{BernoulliRegular.generalizedBernoulliPDivisible_of_nontrivial_oddComponent}
  \leanok
  \uses{thm:stickelberger, def:gen-bernoulli, def:even-odd}
Let $\chi$ be an odd Dirichlet character modulo $p$. If the
$\chi$-eigenspace $\mathrm{Cl}(\mathcal O_{\Q(\zeta_p)})[p]^\chi$ of the
$p$-part of the class group is non-trivial, then
\[
  p \mid \Bgen{1}{\chi}.
\]
\end{theorem}

\begin{proof}\leanok
Project the Stickelberger annihilation relation to the \(\chi\)-eigenspace.
On that eigenspace the Stickelberger element acts by the scalar
\(\Bgen{1}{\chi}\), up to the standard nonzero normalising factor.  If the
\(\chi\)-component of the \(p\)-class group is non-trivial, a scalar that
annihilates it must be divisible by \(p\); otherwise it would be a unit on the
\(\mathbb F_p\)-component and could not kill a nonzero element.  Hence
\(p\mid \Bgen{1}{\chi}\).
\end{proof}
